\scanpage{132}
\mspage{81}

Le grand \unsure{traquenard} \inserted{ici} qu'il faut arriver à exorciser,
c'est qu'il peut arriver que $i\in\mathcal W$, sans qu'on
ait $i'\in\mathcal W$, et aussi $p\in\mathcal W$ sans qu'on ait $q\in\mathcal W$.
Je m'en étais convaincu en 1983 en écrivant
P.S. (\unsure{devant} les \unsure{sommations|miracles} du ``scratchwork''),
et n'ai pas envie de rechercher un \unsure{contrexemple|contre-exemple}
ici, s'il ne se présente de lui-même (ce
qu'il fera sans doute).

Je vais donner un cas important où ces
ennuis disparaissent :

\uline{Proposition} 6 : Considérons les adhérences $\overline U$
et $\overline U'$ de $U$ et $U'$ dans $X$, $X'$ respectivement,
et supposons que \deleted{$U\unsure{,|\to}\overline U$} les inclusions
\[
 U\xrightarrow{i_0}\overline U,\qquad U'\xrightarrow{i'_0}\overline U'
\]
soient \unsure{dans|des} $\mathcal W^{(*)}$. Alors
% slanted
\margin{cf th.~3 p.~47 \unsure{et} pour des conditions générales assurant que cette hypothèse \unsure{est|soit} satisfaite}
\[
\left\{\begin{aligned}
\text{\textcircled{a}}\, i\in\mathcal W&\Longleftrightarrow i'\in\mathcal W &&\text{si on a }\mathrm W(7\ \mathrm{bis})\\
\text{\textcircled{b}}\, p\in\mathcal W&\Longleftrightarrow q\in\mathcal W &&\text{si on a }\mathrm W(7).
\end{aligned}\right.
\]

Démonstration. \deleted{\unsure{La}}Considérons d'abord le cas \deleted{\unsure{a}}\textcircled{a} de
$\unsure{i_0|i}$, $i'$. Soit
\[
 k:\overline U\hookrightarrow X
\]
l'inclusion. Le triangle
\diagram{\begin{tikzcd}[ampersand replacement=\&]
 U \arrow[r,"i_0"] \arrow[rr,bend right=19,"i"'] \& \overline U \arrow[r,"k"] \& X
\end{tikzcd}}
, où $i_0\in\mathcal W$,
montre que $i\in\mathcal W\Longleftrightarrow k\in\mathcal W$. D'autre part $X$ est
la réunion des fermés $\overline U$, $T$, donc on a un
carré cocartésien dans $\mathrm{Cat}$
\diagram{\(\text{(*)}\qquad\begin{tikzcd}[ampersand replacement=\&]
 \overline U\cap T=\dot T \arrow[r,"k_0"] \arrow[d,"j_0"'] \& T \arrow[d,"j"]\\
 \overline U \arrow[r,"k"] \& X
\end{tikzcd}\)}
L'axiome $\mathrm W(7\ \mathrm{bis})$ (sous sa forme duale, (i.e. appliqué
au carré formé par les catégories opposées\ldots) implique

% archive sheet number: 132
